problem:  
Let $(\mathcal{O}, g_{\mathcal{O}})$ be a compact Riemannian $n$‑orbifold with positive sectional curvature. Assume that $\mathcal{O}$ has finitely many *isolated cyclic quotient singularities*, each modeled locally on $\mathbb{R}^n / \Gamma_i$ with $\Gamma_i < \mathrm{SO}(n)$ a finite cyclic group acting freely on $S^{n-1}$. Suppose further that the smooth part $\mathcal{O}_{\mathrm{reg}}$ admits neighborhoods around each singular point whose boundaries are lens spaces $L_i = S^{n-1} / \Gamma_i$ carrying the induced metric of positive sectional curvature.  

Construct a smooth, compact manifold $M$ by removing small neighborhoods of each singular point of $\mathcal{O}$ and gluing in smooth disk bundles $E_{\Gamma_i}$ with $\partial E_{\Gamma_i} = L_i$, producing a natural resolution map $\pi : M \to \mathcal{O}$ that is an isometry on the complement of these neighborhoods.  

**Question:**  
Does there exist a Riemannian metric $g$ on $M$ with positive sectional curvature such that $g = \pi^* g_{\mathcal{O}}$ outside the inserted $E_{\Gamma_i}$, and inside each $E_{\Gamma_i}$ the metric smoothly extends while maintaining positivity of curvature?  
Alternatively, what necessary geometric or topological conditions on the $\Gamma_i$ (or on the topology of the attached $E_{\Gamma_i}$) determine whether such a curvature‑preserving resolution exists?

Why is it a "good" problem:  
This formulation removes the inconsistent assumption of a circle action while preserving the motivating idea: testing the *stability of positive sectional curvature* under the smooth resolution of isolated orbifold singularities. The setting is fully coherent within Riemannian orbifold theory and naturally connected to well‑known local models such as lens spaces arising from cyclic quotients.  

The problem is “good” in multiple senses:  
- **Conceptual Simplicity:** It asks whether positive curvature survives a highly structured, local modification of the manifold.  
- **Interdisciplinary Bridge:** It binds together Riemannian geometry (sectional curvature inequalities), geometric topology (resolution and gluing), and representation theory (actions of $\Gamma_i$ on tangent spaces).  
- **Potential Impact:** Successful resolution methods would yield new examples of positively curved manifolds—objects of extreme rarity—and illuminate how orbifold singularities encode curvature geometrically.  
- **Depth and Difficulty:** Even in low dimensions (e.g., $n=4,5$), no known analytic or geometric framework guarantees curvature preservation under such smoothing. Developing one would require fusing metric gluing techniques, curvature flows, and topological insights.  

Thus, despite its compact statement, this question reaches deeply into the structure of curvature and singularity, epitomizing the “narrow entrance, profound depth” hallmark of a truly research‑level mathematical problem.